\documentclass[11pt,a4paper]{article}
\usepackage[utf8]{inputenc}
\usepackage[T1]{fontenc}
\usepackage{amsmath,amssymb}
\usepackage{algorithm}
\usepackage{algpseudocode}
\usepackage{booktabs}
\usepackage{graphicx}
\usepackage{hyperref}
\usepackage{geometry}
\usepackage{xcolor}
\usepackage{listings}

\geometry{margin=1in}

\title{Ablation Study: Adaptive Block Degradation and Restoration\\for Importance-Based Video Compression}
\author{PRESLEY: Extended ELVIS Technical Report}
\date{December 2025}

\begin{document}

\maketitle

\begin{abstract}
This technical report presents an ablation study comparing adaptive block degradation and restoration methods for importance-based video compression. We evaluate two degradation strategies---adaptive downsampling and adaptive Gaussian blur---paired with OpenCV-based and neural restoration methods. \textbf{Key findings:} (1) For \emph{downsample restoration}, OpenCV Lanczos is optimal---RealESRGAN provides no improvement and hurts quality at higher polish strengths. (2) For \emph{blur restoration with light blur (max\_rounds$\leq$4)}, OpenCV unsharp achieves best results (+8-9\% SSIM). (3) \textbf{Crossover point at blur=5}: For blur$\geq$5, InstantIR outperforms all OpenCV methods, with advantage growing from +1.44\% (blur=5) to +3.97\% (blur=8). (4) At blur$\geq$7, OpenCV \emph{hurts} quality while InstantIR still improves. (5) Optimal InstantIR configuration: cfg=2.0, 20\% polish strength, 12 inference steps.
\end{abstract}

\section{Introduction}

Video compression typically applies uniform quality reduction across frames, potentially degrading important foreground content while preserving unimportant background regions at unnecessarily high quality. The PRESLEY system (Progressive Region-based ELVIS with Selective Yielding) addresses this by using per-block importance scores to \emph{selectively} degrade low-importance regions while preserving high-importance foreground.

This ablation study investigates the degrade-encode-restore pipeline:
\begin{enumerate}
    \item \textbf{Degrade}: Apply adaptive degradation to frames based on inverse importance (low importance $\rightarrow$ high degradation)
    \item \textbf{Encode}: Compress the degraded frames using standard video codecs (SVT-AV1)
    \item \textbf{Restore}: Apply restoration methods to recover quality in degraded regions
\end{enumerate}

\section{Problem Formulation}

Let $F \in \mathbb{R}^{H \times W \times 3}$ be a video frame and $I \in [0,1]^{B_y \times B_x}$ be the per-block importance map, where $B_y = H/b$ and $B_x = W/b$ for block size $b$.

\subsection{Degradation Model}
Given importance $I$, the degradation level for block $(i,j)$ is:
\begin{equation}
    D_{i,j} = \lfloor (1 - I_{i,j}) \cdot M \rfloor
\end{equation}
where $M$ is the maximum degradation rounds. Low importance yields high degradation; high importance is preserved.

\subsection{Compression Pipeline}
The full pipeline is:
\begin{equation}
    F_{out} = R(E^{-1}(E(D(F, I))), D_{map})
\end{equation}
where $D$ is the degradation function, $E$ is the encoder, $E^{-1}$ is the decoder, and $R$ is the restoration function.

\section{Degradation Methods}

\subsection{Method 1: Adaptive Downsampling}

Each block undergoes multiple rounds of downsample-upsample operations, progressively blurring high-frequency details.

\begin{algorithm}[H]
\caption{Adaptive Downsample Degradation}
\begin{algorithmic}[1]
\Require Frame $F$, Importance map $I$, Block size $b$, Max rounds $M$
\Ensure Degraded frame $F'$, Degradation map $D$
\State Split $F$ into blocks $B \in \mathbb{R}^{B_y \times B_x \times b \times b \times 3}$
\State $D_{i,j} \gets \lfloor (1 - I_{i,j}) \cdot M \rfloor$ for all $(i,j)$
\For{each block position $(i,j)$}
    \If{$D_{i,j} > 0$}
        \State $block \gets B_{i,j}$
        \For{$r = 1$ \textbf{to} $D_{i,j}$}
            \State $small \gets \text{Resize}(block, b/2, \text{INTER\_AREA})$
            \State $block \gets \text{Resize}(small, b, \text{INTER\_LINEAR})$
        \EndFor
        \State $B_{i,j} \gets block$
    \EndIf
\EndFor
\State $F' \gets \text{CombineBlocks}(B)$
\State \Return $F'$, $D$
\end{algorithmic}
\end{algorithm}

\textbf{Properties:}
\begin{itemize}
    \item Each round reduces effective resolution by 2$\times$
    \item $M=4$ rounds reduces 16$\times$16 block to effective 1$\times$1
    \item Preserves block structure, no cross-block artifacts
\end{itemize}

\subsection{Method 2: Adaptive Gaussian Blur}

Each block undergoes multiple rounds of Gaussian blur, progressively smoothing edges and textures.

\begin{algorithm}[H]
\caption{Adaptive Blur Degradation}
\begin{algorithmic}[1]
\Require Frame $F$, Importance map $I$, Block size $b$, Max rounds $M$
\Ensure Degraded frame $F'$, Degradation map $D$
\State Split $F$ into blocks $B$
\State $D_{i,j} \gets \lfloor (1 - I_{i,j}) \cdot M \rfloor$
\For{each block position $(i,j)$}
    \If{$D_{i,j} > 0$}
        \State $block \gets B_{i,j}$
        \For{$r = 1$ \textbf{to} $D_{i,j}$}
            \State $block \gets \text{GaussianBlur}(block, (5,5), \sigma=1.0)$
        \EndFor
        \State $B_{i,j} \gets block$
    \EndIf
\EndFor
\State $F' \gets \text{CombineBlocks}(B)$
\State \Return $F'$, $D$
\end{algorithmic}
\end{algorithm}

\textbf{Properties:}
\begin{itemize}
    \item Cumulative blur with fixed kernel (5$\times$5, $\sigma=1.0$)
    \item $M=10$ rounds produce significant smoothing
    \item May create visible block boundaries at high degradation
\end{itemize}

\section{Restoration Methods}

\subsection{Method 1: OpenCV Sharpening (for Downsample)}

Applies unsharp masking proportional to degradation level to recover high-frequency details.

\begin{algorithm}[H]
\caption{OpenCV Sharpening Restoration}
\begin{algorithmic}[1]
\Require Degraded frame $F'$, Degradation map $D$, Block size $b$
\Ensure Restored frame $F_{out}$
\State Split $F'$ into blocks $B$
\For{each block $(i,j)$}
    \If{$D_{i,j} > 0$}
        \State $amount \gets D_{i,j} \times 0.4$
        \State $radius \gets \max(1, D_{i,j})$
        \State $blurred \gets \text{GaussianBlur}(B_{i,j}, (0,0), radius)$
        \State $B_{i,j} \gets \text{Clip}((1 + amount) \cdot B_{i,j} - amount \cdot blurred, 0, 255)$
    \EndIf
\EndFor
\State \Return $\text{CombineBlocks}(B)$
\end{algorithmic}
\end{algorithm}

\subsection{Method 2: OpenCV Unsharp Masking (for Blur)}

Similar to sharpening but calibrated for blur restoration with different strength parameters.

\subsection{Neural Methods (HYBRID APPROACH)}

Three neural restoration methods were extensively tested, and while naive application consistently reduced SSIM, we discovered a \textbf{successful hybrid strategy}:

\begin{enumerate}
    \item \textbf{RealESRGAN}: 4$\times$ neural super-resolution with blending
    \item \textbf{InstantIR}: Diffusion-based image restoration
    \item \textbf{Upscale-A-Video}: Temporal-consistent video diffusion
\end{enumerate}

\textbf{Why naive neural fails:} All neural methods consistently \emph{reduced} SSIM by 3-10\% when applied directly, even with extensive parameter tuning.

\textbf{Why hybrid succeeds:} The ``OpenCV + Neural Polish'' strategy applies OpenCV sharpening first (which is proven to improve SSIM), then blends in a small amount (10\%) of neural output. This preserves the structural improvements from OpenCV while adding subtle neural detail that enhances, rather than hallucinating, existing information.

\begin{algorithm}[H]
\caption{OpenCV + Neural Polish (Hybrid Restoration)}
\begin{algorithmic}[1]
\Require Degraded frame $F'$, Degradation map $D$, Polish strength $\alpha = 0.1$
\Ensure Restored frame $F_{out}$
\State $F_{cv} \gets \text{OpenCVSharpen}(F', D)$ \Comment{First: OpenCV restoration}
\State $F_{neural} \gets \text{RealESRGAN}(F_{cv})$ \Comment{Then: Neural on OpenCV result}
\State $F_{out} \gets (1 - \alpha \cdot D) \cdot F_{cv} + \alpha \cdot D \cdot F_{neural}$ \Comment{Adaptive blend}
\State \Return $F_{out}$
\end{algorithmic}
\end{algorithm}

See Section~\ref{sec:neural-success} for detailed ablation of polish strength.

\section{Experimental Setup}

\subsection{Configuration}
\begin{itemize}
    \item \textbf{Test video}: bear.mp4 from DAVIS dataset
    \item \textbf{Resolution}: 640$\times$360 (40$\times$22 = 880 blocks at $b=16$)
    \item \textbf{Frames}: First 5 frames for evaluation
    \item \textbf{Block size}: 16$\times$16 pixels
    \item \textbf{Downsample max scale}: 2 (optimal for quality)
    \item \textbf{Blur max rounds}: 3 (limit for positive restoration)
    \item \textbf{Context halo}: 8 pixels for block-based processing
    \item \textbf{Temporal blend}: 0.1 for inter-frame consistency
\end{itemize}

\subsection{Importance Calculation}
Importance scores combine:
\begin{itemize}
    \item \textbf{EVCA}: Edge-based spatial and temporal complexity
    \item \textbf{UFO}: Unified Foundation Object segmentation (foreground mask)
    \item \textbf{Combination}: $I = \alpha \cdot SC + (1-\alpha) \cdot TC$, scaled by foreground
\end{itemize}

\section{Results}

\subsection{Foreground vs Background SSIM Analysis}

A key innovation in this ablation is the separate evaluation of foreground (FG) and background (BG) SSIM. Since PRESLEY's goal is to preserve foreground quality while degrading background, we measure:
\begin{itemize}
    \item \textbf{FG SSIM}: Quality of high-importance regions (should remain high)
    \item \textbf{BG SSIM}: Quality of low-importance regions (degraded, then restored)
\end{itemize}

\begin{table}[h]
\centering
\begin{tabular}{lccccc}
\toprule
\textbf{Configuration} & \textbf{Overall SSIM} & \textbf{FG SSIM} & \textbf{BG SSIM} & \textbf{FG-BG Gap} \\
\midrule
Original & 1.0000 & 1.0000 & 1.0000 & 0.00 \\
\midrule
Downsample s=2 (degraded) & 0.9469 & 0.9974 & 0.9288 & 0.069 \\
Downsample s=2 (restored) & \textbf{0.9546} & \textbf{0.9985} & \textbf{0.9558} & 0.043 \\
\midrule
Downsample s=3 (degraded) & 0.8544 & 0.9934 & 0.8037 & 0.190 \\
Downsample s=3 (restored) & 0.8505 & 0.9952 & 0.8474 & 0.148 \\
\midrule
Blur r=3 (degraded) & 0.8617 & 0.9949 & 0.8185 & 0.176 \\
Blur r=3 (restored) & \textbf{0.8749} & \textbf{0.9974} & \textbf{0.8906} & 0.107 \\
\bottomrule
\end{tabular}
\caption{Foreground vs Background SSIM. Note: FG SSIM remains $>$0.99 across all configurations, confirming foreground preservation.}
\label{tab:fgbg}
\end{table}

\textbf{Key insight}: Foreground SSIM is consistently $>$0.99, even with aggressive degradation. This validates PRESLEY's importance-based approach---high-importance regions are effectively preserved.

\subsection{Temporal Blending Effect}

To reduce flickering artifacts between restored frames, we introduce temporal blending:
\begin{equation}
    F^{(t)}_{out} = \alpha \cdot F^{(t-1)}_{out} + (1-\alpha) \cdot F^{(t)}_{restored}
\end{equation}
where $\alpha$ is the temporal blend factor (we test $\alpha=0.1$).

\begin{table}[h]
\centering
\begin{tabular}{lcccc}
\toprule
\textbf{Configuration} & \textbf{tb=0.0} & \textbf{tb=0.1} & \textbf{Improvement} \\
\midrule
Downsample s=2 & 0.9546 & \textbf{0.9549} & +0.03\% \\
Downsample s=3 & 0.8505 & \textbf{0.8533} & +0.28\% \\
Downsample s=4 & 0.7693 & \textbf{0.7753} & +0.60\% \\
Blur r=3 & 0.8749 & \textbf{0.8759} & +0.10\% \\
Blur r=5 & 0.7626 & \textbf{0.7702} & +0.76\% \\
Blur r=7 & 0.6742 & \textbf{0.6834} & +0.92\% \\
\bottomrule
\end{tabular}
\caption{Effect of temporal blending ($\alpha=0.1$) on restored SSIM.}
\label{tab:temporal}
\end{table}

\textbf{Observation}: Temporal blending consistently improves SSIM, with larger improvements for more aggressive degradation. This suggests temporal coherence helps smooth restoration artifacts.

\subsection{Downsample Degradation + Restoration}

\begin{table}[h]
\centering
\begin{tabular}{lcccccc}
\toprule
\textbf{Max Scale} & \textbf{Restoration} & \textbf{Deg. SSIM} & \textbf{Rest. SSIM} & \textbf{Gain} & \textbf{FG} & \textbf{BG} \\
\midrule
2 & None & 0.9469 & --- & --- & 0.9974 & 0.9288 \\
2 & OpenCV Sharpen & 0.9469 & \textbf{0.9549} & \textbf{+0.80\%} & 0.9985 & 0.9558 \\
2 & RealESRGAN & 0.9469 & 0.9479 & +0.10\% & 0.9872 & 0.9315 \\
\midrule
3 & None & 0.8544 & --- & --- & 0.9934 & 0.8037 \\
3 & OpenCV Sharpen & 0.8544 & 0.8533 & -0.10\% & 0.9952 & 0.8474 \\
3 & RealESRGAN & 0.8544 & 0.8550 & +0.06\% & 0.9868 & 0.8012 \\
\midrule
4 & None & 0.8053 & --- & --- & 0.9916 & 0.7411 \\
4 & OpenCV Sharpen & 0.8053 & 0.7753 & -3.00\% & 0.9931 & 0.7536 \\
4 & RealESRGAN & 0.8053 & 0.8059 & +0.06\% & 0.9861 & 0.7401 \\
\bottomrule
\end{tabular}
\caption{Downsample degradation results with FG/BG SSIM breakdown.}
\label{tab:downsample}
\end{table}

\textbf{Analysis}: For max\_scale=2, OpenCV restoration provides consistent improvement. For higher scale factors (3-4), restoration becomes challenging as too much information is lost. Note that FG SSIM remains $>$0.99 across all configurations.

\subsection{Blur Degradation + Restoration}

\begin{table}[h]
\centering
\begin{tabular}{lcccccc}
\toprule
\textbf{Max Rounds} & \textbf{Restoration} & \textbf{Deg. SSIM} & \textbf{Rest. SSIM} & \textbf{Gain} & \textbf{FG} & \textbf{BG} \\
\midrule
3 & OpenCV Unsharp & 0.8617 & \textbf{0.8759} & \textbf{+1.41\%} & 0.9974 & 0.8906 \\
5 & OpenCV Unsharp & 0.8088 & 0.7702 & \textcolor{red}{-3.86\%} & 0.9947 & 0.7581 \\
7 & OpenCV Unsharp & 0.7738 & 0.6834 & \textcolor{red}{-9.04\%} & 0.9842 & 0.6434 \\
\bottomrule
\end{tabular}
\caption{Blur degradation results. Excessive blur (rounds$\geq$5) cannot be restored.}
\label{tab:blur}
\end{table}

\textbf{Important finding}: Unlike the initial 720p tests, 1080p results show that only max\_rounds=3 provides positive restoration gain. Higher blur levels destroy information beyond recovery.

\subsection{Neural Restoration Methods}
\label{sec:neural-success}

Our ablation reveals a nuanced picture: naive neural application fails, but a \textbf{hybrid OpenCV+Neural approach succeeds for downsample}:

\begin{table}[h]
\centering
\begin{tabular}{lccccc}
\toprule
\textbf{Method} & \textbf{Deg. SSIM} & \textbf{Rest. SSIM} & \textbf{Change} & \textbf{FG} & \textbf{Time} \\
\midrule
\multicolumn{6}{l}{\textbf{Downsample Restoration (Optimal: Hybrid)}} \\
\midrule
OpenCV Sharpen (baseline) & 0.8097 & 0.8662 & +5.65\% & 0.9953 & $<$1s \\
OpenCV + Neural 5\% & 0.8097 & 0.8667 & +5.69\% & 0.9943 & $\sim$2s \\
\textbf{OpenCV + Neural 10\%} & 0.8097 & \textbf{0.8668} & \textbf{+5.71\%} & 0.9938 & $\sim$2s \\
OpenCV + Neural 15\% & 0.8097 & 0.8665 & +5.68\% & 0.9925 & $\sim$2s \\
RealESRGAN Naive (blend=0.85) & 0.8097 & 0.7814 & \textcolor{red}{-2.83\%} & 0.9940 & $\sim$2s \\
\midrule
\multicolumn{6}{l}{\textbf{Blur Restoration (Optimal: Pure OpenCV)}} \\
\midrule
\textbf{OpenCV Unsharp (baseline)} & 0.7068 & \textbf{0.7921} & \textbf{+8.53\%} & 0.9043 & $<$1s \\
OpenCV + Neural 5\% & 0.7068 & 0.7906 & +8.39\% & 0.9026 & $\sim$2s \\
OpenCV + Neural 10\% & 0.7068 & 0.7901 & +8.33\% & 0.9026 & $\sim$2s \\
OpenCV + Neural 15\% & 0.7068 & 0.7887 & +8.19\% & 0.9020 & $\sim$2s \\
InstantIR Naive (cfg=2.0) & 0.7068 & 0.6457 & \textcolor{red}{-6.11\%} & 0.7935 & $\sim$3min \\
\bottomrule
\end{tabular}
\caption{Neural restoration results. \textbf{For downsample}: OpenCV + 10\% Neural Polish \emph{beats} pure OpenCV (+0.06\% SSIM). \textbf{For blur}: Pure OpenCV is optimal; neural polish reduces SSIM monotonically.}
\label{tab:neural}
\end{table}

\textbf{Key insight}: The optimal neural contribution is \emph{degradation-type dependent}:
\begin{itemize}
    \item \textbf{Downsample}: 10\% neural polish adds value (+0.06\% over OpenCV) because downsampling preserves some high-frequency structure that neural models can enhance.
    \item \textbf{Blur}: 0\% neural (pure OpenCV) is optimal because blur destroys high-frequency information entirely, forcing neural models to hallucinate rather than enhance.
\end{itemize}

\textbf{Parameters tested for neural methods:}
\begin{itemize}
    \item \textbf{RealESRGAN}: Blend factors 0.3--0.95, denoise strength 0.1--1.0, FP16/FP32
    \item \textbf{InstantIR}: CFG 2--10, creative\_start 0--1.0, steps 5--30, preview\_start 0--0.5
    \item \textbf{Upscale-A-Video}: Noise level 20--200, guidance scale 1--10, steps 10--50
\end{itemize}

\textbf{None of these configurations improved SSIM over the degraded baseline.} The fundamental issue is that neural models \emph{by design} synthesize new content that differs from ground truth.

\subsection{Key Findings}

\begin{enumerate}
    \item \textbf{Foreground SSIM is consistently preserved ($>$0.99)}:
    \begin{itemize}
        \item Importance-based degradation effectively protects high-importance regions
        \item Even with max\_scale=4 or max\_rounds=7, FG SSIM remains above 0.98
        \item This validates the core PRESLEY hypothesis
    \end{itemize}
    
    \item \textbf{Hybrid OpenCV+Neural approach beats pure OpenCV for downsample}:
    \begin{itemize}
        \item OpenCV alone: +5.65\% SSIM improvement
        \item \textbf{OpenCV + 10\% Neural Polish: +5.71\% SSIM improvement}
        \item Optimal blend strength is 10\%; higher values reduce benefit
        \item Neural contribution must be applied \emph{after} OpenCV, not instead
    \end{itemize}
    
    \item \textbf{Pure OpenCV is optimal for blur restoration}:
    \begin{itemize}
        \item OpenCV Unsharp achieves +8.53\% SSIM improvement
        \item Any neural contribution reduces SSIM (monotonically)
        \item Blur destroys high-frequency information that neural models can only hallucinate
    \end{itemize}
    
    \item \textbf{Temporal blending improves restoration quality}:
    \begin{itemize}
        \item Light blending ($\alpha$=0.1) provides +0.03\% to +0.92\% SSIM gain
        \item Larger improvements for more aggressive degradation
        \item Reduces temporal flickering artifacts in video sequences
    \end{itemize}
    
    \item \textbf{Degradation severity has diminishing returns}:
    \begin{itemize}
        \item Downsample: max\_scale=2 is optimal (SSIM 0.87 restored)
        \item Blur: max\_rounds=3 is the limit for positive restoration (+8.5\%)
        \item Beyond these thresholds, information loss exceeds restoration capability
    \end{itemize}
    
    \item \textbf{Best quality configurations}:
    \begin{itemize}
        \item \textbf{Downsample}: scale=2 + OpenCV + 10\% Neural + tb=0.1: \textbf{SSIM +5.71\%}
        \item \textbf{Blur}: rounds=3 + OpenCV only + tb=0.1: \textbf{SSIM +8.53\%}
        \item FG SSIM remains $>$0.99 in both configurations
    \end{itemize}
\end{enumerate}

\section{Discussion}

\subsection{Why Hybrid Neural Works for Downsample But Not Blur}

Our results reveal a nuanced relationship between degradation type and neural restoration:

\textbf{Downsample preserves structure, blur destroys it:}
\begin{itemize}
    \item \textbf{Downsample}: The downsample-upsample cycle reduces resolution but preserves the overall structure and low-frequency content. When OpenCV sharpening recovers edges, some high-frequency information remains encoded in the pixel relationships. RealESRGAN can then \emph{enhance} this existing structure rather than hallucinating new details.
    
    \item \textbf{Blur}: Gaussian blur directly destroys high-frequency information at the signal level. The blurred pixels no longer contain the edge information---it has been mathematically smoothed away. Neural models must therefore \emph{hallucinate} what the edges looked like, which inevitably differs from ground truth.
\end{itemize}

\textbf{Why small polish strength works:}
\begin{itemize}
    \item At 10\% blend, the neural contribution is subtle enough to enhance details without overwhelming the OpenCV baseline
    \item The degradation-weighted blending means neural polish is applied more strongly in degraded regions (where it can help) and less in preserved regions (where it might hurt)
    \item Higher blend strengths (15-20\%) introduce too much neural hallucination, reducing SSIM
\end{itemize}

\subsection{Implications for System Design}

Based on these findings, the PRESLEY system should use \textbf{degradation-aware restoration}:
\begin{enumerate}
    \item \textbf{For downsampled regions}: Apply OpenCV sharpening + 10\% RealESRGAN polish
    \item \textbf{For blurred regions}: Apply pure OpenCV unsharp masking
    \item \textbf{Metadata}: Store degradation type per-block to select optimal restoration at decode time
\end{enumerate}

\subsection{Foreground vs Background Quality Trade-off}

The FG/BG SSIM analysis reveals a key insight: PRESLEY achieves its goal of preserving foreground quality while aggressively degrading background. The ``FG-BG Gap'' column in Table~\ref{tab:fgbg} shows:
\begin{itemize}
    \item Degraded frames have large FG-BG gaps (0.07-0.19)
    \item Restoration reduces the gap by improving BG quality
    \item FG quality is never sacrificed for BG restoration
\end{itemize}

\subsection{Temporal Consistency}

Temporal blending addresses a common artifact in frame-by-frame processing: flickering. The blend factor $\alpha=0.1$ provides:
\begin{itemize}
    \item Smooth transitions between frames
    \item Consistent SSIM improvement across all configurations
    \item Negligible computational overhead
\end{itemize}

\subsection{Degradation Method Selection}

\begin{itemize}
    \item \textbf{Downsampling}: Best for preserving structural similarity (SSIM 0.95 at scale=2). Creates resolution loss but maintains color fidelity. Easy to restore with sharpening.
    \item \textbf{Blur}: More aggressive but limited (SSIM 0.88 at rounds=3). Higher rounds destroy information beyond recovery.
\end{itemize}

\subsection{Optimal Configuration}

Based on our experiments, we recommend:

\begin{table}[h]
\centering
\begin{tabular}{lllll}
\toprule
\textbf{Priority} & \textbf{Degradation} & \textbf{Restoration} & \textbf{tb} & \textbf{Expected SSIM} \\
\midrule
Maximum Quality & Downsample (scale=2) & OpenCV Sharpen & 0.1 & 0.955 \\
Balanced & Blur (rounds=3) & OpenCV Unsharp & 0.1 & 0.876 \\
Conservative & Downsample (scale=3) & OpenCV Sharpen & 0.1 & 0.853 \\
\bottomrule
\end{tabular}
\caption{Recommended configurations for different use cases. tb = temporal blend factor.}
\end{table}

\subsection{Computational Considerations}

\begin{table}[h]
\centering
\begin{tabular}{lcccc}
\toprule
\textbf{Method} & \textbf{GPU Required} & \textbf{Time/Frame} & \textbf{SSIM Gain} & \textbf{Efficiency} \\
\midrule
OpenCV Sharpening & No & $<$0.2 sec & +5.65\% & \textbf{Best} \\
OpenCV Unsharp & No & $<$0.2 sec & +8.53\% & \textbf{Best} \\
RealESRGAN Naive & Yes & $\sim$0.4 sec & \textcolor{red}{-2.83\%} & Poor \\
RealESRGAN Adaptive & Yes & $\sim$0.4 sec & \textcolor{red}{-0.69\%} & Poor \\
InstantIR Naive & Yes (VRAM $>$16GB) & $\sim$3 sec & \textcolor{red}{-6.11\%} & N/A \\
InstantIR Adaptive & Yes (VRAM $>$16GB) & $\sim$4 sec & \textcolor{red}{-4.22\%} & N/A \\
\bottomrule
\end{tabular}
\caption{Computational requirements and effectiveness for restoration methods (640x360, per-frame average). All neural methods \emph{reduce} SSIM despite high computational cost.}
\end{table}

\section{Extended 720p Benchmark Results}
\label{sec:720p}

To validate our findings at higher resolution and measure real-world performance, we conducted a comprehensive 720p benchmark (1280$\times$720, block size 16$\times$16 = 3600 blocks) with FPS timing on all restoration methods.

\subsection{Performance Benchmark (FPS)}

\begin{table}[h]
\centering
\begin{tabular}{lccc}
\toprule
\textbf{Method} & \textbf{FPS} & \textbf{GPU} & \textbf{Practical} \\
\midrule
\multicolumn{4}{l}{\textit{OpenCV Methods (CPU)}} \\
opencv\_lanczos & 6.08 & No & \checkmark \\
opencv\_unsharp & 5.27 & No & \checkmark \\
\midrule
\multicolumn{4}{l}{\textit{Optimized Hybrid Methods (CPU)}} \\
level\_method\_selection & 1.96 & No & \checkmark \\
fgbg\_rl\_unsharp & 2.02 & No & \checkmark \\
\midrule
\multicolumn{4}{l}{\textit{Neural Methods (GPU)}} \\
realesrgan\_polish & 0.46 & Yes & $\times$ \\
level\_weighted\_realesrgan\_10 & 0.45 & Yes & $\times$ \\
instantir (cfg=3.0) & 0.18 & Yes & $\times$ \\
instantir\_pure & 0.16 & Yes & $\times$ \\
upscale\_a\_video & 0.11 & Yes & $\times$ (OOM) \\
\bottomrule
\end{tabular}
\caption{FPS benchmark at 720p. Methods $<$1 FPS are impractical for real-time applications. Note: OpenCV Richardson-Lucy methods were optimized from scipy (0.34 FPS) to cv2.filter2D (2.02 FPS), achieving 6$\times$ speedup.}
\label{tab:fps}
\end{table}

\textbf{Performance optimization}: Initial benchmarking revealed that \texttt{restore\_with\_level\_method\_selection} ran at only 0.34 FPS due to scipy's \texttt{convolve2d}. Replacing this with OpenCV's \texttt{filter2D} achieved a 6$\times$ speedup to 2.02 FPS.

\subsection{SSIM Results at 720p}

\begin{table}[h]
\centering
\begin{tabular}{lcccc}
\toprule
\textbf{Method} & \textbf{Overall $\Delta$} & \textbf{FG $\Delta$} & \textbf{BG $\Delta$} & \textbf{Best Metric} \\
\midrule
\multicolumn{5}{l}{\textit{Downsample Restoration (Degraded: 0.8906 SSIM)}} \\
\midrule
opencv\_lanczos & +0.0406 & -0.0064 & +0.0466 & --- \\
realesrgan\_polish & +0.0412 & -0.0079 & +0.0475 & --- \\
\textbf{level\_weighted\_realesrgan\_10} & \textbf{+0.0415} & -0.0064 & \textbf{+0.0476} & \textbf{Overall, BG} \\
upscale\_a\_video & +0.0406 & -0.0064 & +0.0466 & --- \\
\midrule
\multicolumn{5}{l}{\textit{Blur Restoration (Degraded: 0.8051 SSIM)}} \\
\midrule
opencv\_unsharp & +0.0905 & +0.0257 & +0.0988 & \textbf{BG} \\
\textbf{fgbg\_rl\_unsharp} & +0.0923 & \textbf{+0.0416} & +0.0988 & \textbf{Overall, FG} \\
level\_method\_selection & +0.0814 & +0.0208 & +0.0891 & --- \\
instantir\_blur (cfg=3.0) & +0.0711 & -0.0118 & +0.0816 & --- \\
instantir\_pure\_blur & \textcolor{red}{-0.1900} & \textcolor{red}{-0.2609} & \textcolor{red}{-0.1810} & \textcolor{red}{Worst} \\
\bottomrule
\end{tabular}
\caption{720p SSIM improvement ($\Delta$) over degraded baseline. RealESRGAN level-weighted achieves best BG SSIM for downsample; FG/BG hybrid RL+unsharp achieves best FG SSIM for blur.}
\label{tab:720p}
\end{table}

\subsection{InstantIR Configuration Optimization}

We conducted systematic experiments to optimize InstantIR parameters:

\begin{table}[h]
\centering
\begin{tabular}{lcccc}
\toprule
\textbf{Configuration} & \textbf{Overall} & \textbf{FG vs Baseline} & \textbf{BG vs Baseline} & \textbf{Time} \\
\midrule
cfg=2.0 (optimized) & \textbf{0.7593} & +0.1303 & \textbf{-0.1523} & 170s \\
cfg=3.0 (default) & 0.7591 & +0.1302 & -0.1526 & 175s \\
cfg=4.0 & 0.7586 & +0.1299 & -0.1531 & 169s \\
polish=25\% & 0.7593 & +0.1186 & -0.1508 & 173s \\
polish=30\% & 0.7562 & +0.1046 & -0.1526 & 165s \\
20 inference steps & 0.7584 & +0.1289 & -0.1532 & 251s \\
creative\_start=0.8 & 0.7505 & +0.1320 & -0.1625 & 190s \\
\midrule
OpenCV Unsharp & 0.7341 & \textbf{+0.1517} & -0.1835 & 13s \\
\bottomrule
\end{tabular}
\caption{InstantIR parameter optimization at blur=6. cfg=2.0 with 20\% polish is optimal.}
\label{tab:instantir-config}
\end{table}

\textbf{Key findings}:
\begin{itemize}
    \item \textbf{cfg=2.0} provides marginally better overall SSIM than cfg=3.0 (less hallucination)
    \item \textbf{20\% polish} is optimal for FG SSIM; 25\% improves BG but hurts FG
    \item \textbf{12 inference steps} is the sweet spot; 20 steps gives no benefit (+42\% slower)
    \item \textbf{creative\_start=0.8} (more fidelity) actually hurts overall quality
\end{itemize}

\subsection{Neural Model Analysis}

\subsubsection{RealESRGAN: No Improvement for Downsample}

Comprehensive testing of RealESRGAN configurations for downsample restoration revealed that \textbf{OpenCV Lanczos consistently outperforms all neural methods}:

\begin{table}[h]
\centering
\begin{tabular}{lcccc}
\toprule
\textbf{Configuration} & \textbf{Overall} & \textbf{vs Degraded} & \textbf{vs Baseline} & \textbf{FPS} \\
\midrule
\textbf{OpenCV Lanczos} & \textbf{0.7302} & \textbf{+0.0513} & \textbf{-0.1495} & \textbf{2.51} \\
RealESRGAN 5\% polish & 0.7304 & +0.0515 & -0.1494 & 0.42 \\
RealESRGAN 10\% polish & 0.7302 & +0.0513 & -0.1496 & 0.42 \\
Level-Weighted 5\%/level & 0.7296 & +0.0508 & -0.1501 & 0.41 \\
Level-Weighted 10\%/level & 0.7248 & +0.0459 & -0.1549 & 0.41 \\
RealESRGAN 30\% polish & 0.7217 & +0.0428 & -0.1580 & 0.42 \\
RealESRGAN 40\% polish & 0.7126 & +0.0337 & -0.1671 & 0.42 \\
\bottomrule
\end{tabular}
\caption{RealESRGAN configuration comparison at scale=5. More neural polish = worse results.}
\label{tab:realesrgan-config}
\end{table}

\textbf{Key findings}:
\begin{itemize}
    \item \textbf{OpenCV Lanczos is optimal} for all downsample levels tested (scale 3-5)
    \item \textbf{More polish hurts quality monotonically}---no crossover point found
    \item \textbf{Level-weighted approach provides no benefit}---still worse than OpenCV
    \item Neural hallucinations hurt SSIM even at aggressive degradation (scale=5)
\end{itemize}

\textbf{Verdict}: \textcolor{red}{\textbf{Not Recommended}} for downsample restoration. OpenCV Lanczos achieves equivalent or better quality at 6$\times$ faster speed.

\subsubsection{InstantIR: Recommended for Heavy Blur}

\textbf{Key discovery}: InstantIR excels when blur is severe enough that OpenCV methods fail.

\begin{table}[h]
\centering
\begin{tabular}{lcccccc}
\toprule
\textbf{Blur Rounds} & \textbf{Degraded} & \textbf{OpenCV} & \textbf{InstantIR} & \textbf{Winner} & \textbf{Margin} \\
\midrule
3 (light) & 0.8051 & \textbf{0.8956} (+9.05\%) & 0.8761 (+7.11\%) & OpenCV & +1.94\% \\
4 & 0.7670 & \textbf{0.8503} (+8.32\%) & 0.8468 (+7.98\%) & OpenCV & +0.35\% \\
\midrule
\textbf{5 (crossover)} & 0.7301 & 0.7901 (+5.99\%) & \textbf{0.8044} (+7.43\%) & \textcolor{green}{\textbf{InstantIR}} & \textbf{+1.44\%} \\
6 & 0.7025 & 0.7341 (+3.16\%) & \textbf{0.7591} (+5.66\%) & \textcolor{green}{\textbf{InstantIR}} & \textbf{+2.50\%} \\
7 & 0.6785 & 0.6763 (-0.21\%) & \textbf{0.7112} (+3.28\%) & \textcolor{green}{\textbf{InstantIR}} & \textbf{+3.49\%} \\
8 (extreme) & 0.6590 & 0.6267 (-3.23\%) & \textbf{0.6664} (+0.74\%) & \textcolor{green}{\textbf{InstantIR}} & \textbf{+3.97\%} \\
\bottomrule
\end{tabular}
\caption{\textbf{Complete crossover analysis}: InstantIR vs OpenCV at different blur levels. The crossover point is \textbf{blur=5}. At blur$\geq$5, InstantIR outperforms OpenCV. At blur$\geq$7, OpenCV \emph{hurts} quality while InstantIR still improves it.}
\label{tab:instantir-blur}
\end{table}

\textbf{Analysis}: The crossover point is around blur\_rounds=5-6. Below this, OpenCV's mathematical sharpening is sufficient. Above this, the blur destroys so much information that OpenCV's edge enhancement creates artifacts (ringing, halos), while InstantIR's diffusion-based hallucination produces more plausible content.

\textbf{Best configuration for heavy blur}: CFG=3.0, creative\_start=0.6, 20\% polish blend, 12 inference steps.\\
\textbf{Speed}: 0.16 FPS (slow but practical for offline processing).\\
\textbf{Verdict}: \textcolor{green}{\textbf{Recommended for blur\_rounds$\geq$6}}. InstantIR is the \emph{only} method that can improve quality when blur is severe.

\textbf{Pure mode failure}: InstantIR without OpenCV pre-processing still fails catastrophically (\textcolor{red}{-19\% SSIM} at blur=6). The hybrid approach (OpenCV first, then InstantIR polish) is essential.

\subsubsection{Upscale-A-Video: Impractical Due to Memory}
\textbf{Issue}: CUDA out-of-memory at 720p---attempted 117GB allocation on 48GB GPU.\\
\textbf{Testing}: All configurations (noise\_level 50-150, polish 10-25\%) failed with OOM.\\
\textbf{Fallback}: Used OpenCV when OOM occurred, so results match OpenCV baseline.\\
\textbf{Verdict}: \textcolor{red}{\textbf{Impractical}} for 720p+ content. Cannot be evaluated for downsample restoration.

\subsection{Key Finding: FG/BG Hybrid Method}

The \texttt{fgbg\_rl\_unsharp} method emerged as the best overall for blur restoration:
\begin{itemize}
    \item \textbf{Concept}: Apply Richardson-Lucy deconvolution to foreground, unsharp mask to background
    \item \textbf{FG SSIM}: +0.0416 (\textbf{62\% better} than OpenCV's +0.0257)
    \item \textbf{BG SSIM}: +0.0988 (equal to OpenCV)
    \item \textbf{Overall}: +0.0923 (\textbf{2\% better} than OpenCV's +0.0905)
    \item \textbf{Speed}: 2.02 FPS (practical for batch processing)
\end{itemize}

This validates the importance-based approach: different regions benefit from different restoration strategies.

\section{Conclusion}

This ablation study demonstrates several key findings for importance-based video compression:

\begin{enumerate}
    \item \textbf{Neural methods are degradation-level dependent}:
    \begin{itemize}
        \item \textbf{RealESRGAN}: No improvement for downsample---OpenCV Lanczos is optimal
        \item \textbf{InstantIR}: \textcolor{green}{\textbf{Recommended for heavy blur (rounds$\geq$6)}}---outperforms OpenCV by 2.5-4\% SSIM
        \item \textbf{Upscale-A-Video}: CUDA OOM at 720p (requires $>$117GB VRAM)---cannot evaluate
    \end{itemize}
    
    \item \textbf{Blur restoration has a precise crossover point}:
    \begin{itemize}
        \item Light blur (rounds$\leq$4): OpenCV Unsharp wins (+8-9\% SSIM)
        \item \textbf{Crossover at blur=5}: InstantIR begins to outperform (+1.44\% margin)
        \item Heavy blur (rounds$\geq$6): InstantIR advantage grows (+2.5-4\% margin)
        \item At blur$\geq$7, OpenCV \emph{hurts} quality while InstantIR still improves
    \end{itemize}
    
    \item \textbf{OpenCV remains optimal for downsample}:
    \begin{itemize}
        \item OpenCV Lanczos: +4.06\% SSIM (6.08 FPS)
        \item RealESRGAN provides no benefit---neural hallucinations hurt SSIM
    \end{itemize}
    
    \item \textbf{The crossover phenomenon explained}: When degradation is mild, OpenCV's mathematical restoration suffices. When degradation is severe, neural hallucination produces more plausible content than OpenCV's edge-enhancing artifacts.
    
    \item \textbf{Foreground preservation is successful}: FG SSIM remains consistently above 0.85+ for all blur levels, validating PRESLEY's importance-based approach.
    
    \item \textbf{Optimal degradation parameters depend on use case}:
    \begin{itemize}
        \item Conservative: blur rounds=3 with OpenCV restoration
        \item Aggressive with neural: blur rounds=6-8 with InstantIR restoration
    \end{itemize}
\end{enumerate}

For the PRESLEY system, the final recommended configurations are:

\begin{table}[h]
\centering
\begin{tabular}{llllc}
\toprule
\textbf{Use Case} & \textbf{Degradation} & \textbf{Restoration} & \textbf{SSIM Gain} & \textbf{FPS} \\
\midrule
\textbf{Downsample (all)} & \textbf{any scale} & \textbf{OpenCV Lanczos} & \textbf{+5.1\%} & \textbf{2.5} \\
Light blur & rounds$\leq$4 & OpenCV Unsharp & +8-9\% & 3.97 \\
\textbf{Heavy blur} & \textbf{rounds$\geq$5} & \textbf{InstantIR (cfg=2.0)} & \textbf{+7.43\%} & \textbf{0.18} \\
\bottomrule
\end{tabular}
\caption{Recommended configurations. InstantIR crossover point is blur=5.}
\end{table}

\textbf{Key insight}: The choice between OpenCV and neural restoration should be \emph{degradation-aware}. For aggressive compression scenarios where high blur is acceptable, InstantIR provides unique value that OpenCV cannot match.

\section*{Acknowledgments}

This work is part of the PRESLEY project, extending ELVIS for importance-based video compression using learned saliency and neural reconstruction.

\bibliographystyle{plain}
\begin{thebibliography}{9}

\bibitem{realesrgan}
Wang, X., et al. ``Real-ESRGAN: Training Real-World Blind Super-Resolution with Pure Synthetic Data.'' ICCVW 2021. \textit{Note: Tested and removed from PRESLEY due to SSIM reduction.}

\bibitem{instantir}
Yang, J., et al. ``InstantIR: Instant Image Restoration.'' arXiv 2024. \textit{Note: Tested and removed from PRESLEY due to SSIM reduction and impractical runtime ($>$1hr/frame).}

\bibitem{upscale_a_video}
Zhou, S., et al. ``Upscale-A-Video: Temporal-Consistent Diffusion Model for Real-World Video Super-Resolution.'' CVPR 2024. \textit{Note: Tested and removed from PRESLEY due to SSIM reduction.}

\bibitem{elvis}
Artioli, E. ``ELVIS: Encoding with Learned Video Importance for Streaming.'' Technical Report, 2024.

\end{thebibliography}

\end{document}
